% Required package: \usepackage{booktabs}
\subsection{Baseline NPU Architecture and Simulation Flow}
\label{sec:onnxim-architecture}

We evaluate the proposed mechanism in ONNXim using the weight-stationary
configuration summarized in Table~\ref{tab:baseline-configuration}. The
simulated NPU contains four homogeneous cores clocked at 1~GHz. Each core has
a $128\times128$ weight-stationary systolic array, a main scratchpad, and an
accumulator scratchpad. Both scratchpads are double buffered. The configured
main and accumulator capacities are 32~MiB and 4~MiB per core, respectively;
consequently, each active buffer receives one half of the corresponding
capacity. Operands use 2-byte precision, and memory accesses are divided into
32-byte requests.

\begin{table}[t]
  \centering
  \caption{Evaluated ONNXim configuration.}
  \label{tab:baseline-configuration}
  \begin{tabular}{ll}
    \toprule
    Parameter & Value \\
    \midrule
    Number of cores & 4 \\
    Core organization & Weight stationary \\
    SA dimensions per core & $128\times128$ \\
    Core frequency & 1~GHz \\
    Main scratchpad per core & 32~MiB, double buffered \\
    Accumulator scratchpad per core & 4~MiB, double buffered \\
    Operand precision & 2~bytes \\
    Memory-request granularity & 32~bytes \\
    Scheduler & Simple \\
    Interconnect model & BookSim2, 1~GHz \\
    Injection ports per core & 16 \\
    DRAM model & Ramulator2 HBM2, 1.2~GHz \\
    DRAM channels & 16 \\
    Tensor layout & NHWC \\
    \bottomrule
  \end{tabular}
\end{table}

\subsubsection{Operation Tiling and Instruction Generation}

ONNXim first lowers each supported model operation into one or more tiles
according to the operation shape, the selected mapping, the $128\times128$
array dimensions, and scratchpad-capacity constraints. A tile contains an
ordered deque of instructions rather than a single atomic operation. The
instruction set includes memory transfers (\texttt{MOVIN} and
\texttt{MOVOUT}), SA operations (\texttt{GEMM\_PRELOAD} and
\texttt{GEMM}), and non-SA operations executed by the remaining modeled
compute path.

For weight-stationary matrix multiplication, instruction generation first
creates the required \texttt{MOVIN} instructions for weights, activations,
and, when present, bias data. It then emits SA instructions followed by the
required output transfer. A \texttt{GEMM\_PRELOAD} instruction loads a new
weight tile into the array and performs the associated matrix multiplication;
a subsequent \texttt{GEMM} can reuse the resident weights. Thus, only
\texttt{GEMM\_PRELOAD} and \texttt{GEMM} intervals are classified as SA
execution in this study.

\subsubsection{Tile Scheduling and Core Admission}

The evaluated configuration uses ONNXim's simple scheduler. When a layer
becomes executable, its tiles enter the scheduler's global executable queue.
Tiles without a preassigned core are distributed across the four cores in
round-robin order and placed in a per-core executable-tile queue. Layer
barriers prevent successor work from passing a layer whose launched tiles
have not finished.

On every core cycle, the simulator first returns completed tiles to the
scheduler and then attempts to admit a new tile. A core accepts a tile only
when it has fewer than two resident tiles; this is the implemented tile-level
double-buffering limit. Admission selects the front tile of the corresponding
per-core scheduler queue. On admission, the core assigns the tile to the
alternate scratchpad and accumulator-buffer identifiers and flushes the newly
selected main-scratchpad buffer. The baseline scheduler, admission order, and
tile mapping are unchanged by our characterization instrumentation.

\subsubsection{Core Queues and Instruction Flow}

Each resident tile retains its ordered instruction deque. During a core
cycle, the core examines resident tiles in round-robin order and issues at
most one instruction successfully. A successfully issued instruction is
removed from its tile and routed according to its opcode:
\begin{itemize}
  \item \texttt{MOVIN} enters the load-instruction queue;
  \item \texttt{MOVOUT} enters the store-instruction queue; and
  \item an execution instruction enters the single-entry execution queue
        after its source operands become valid.
\end{itemize}
The execution queue subsequently dispatches SA instructions to the SA compute
pipeline and other supported compute instructions to their modeled execution
path. The core marks a tile complete only after its instruction deque is empty
and its final issued instruction has completed.

A load instruction allocates an entry in the selected scratchpad buffer and
is decomposed into 32-byte memory requests. These requests enter the core's
FIFO request queue and then traverse BookSim2 to the appropriate HBM2 channel
modeled by Ramulator2. Read requests remain in an outstanding-access table
after leaving the core FIFO. Returned responses traverse the interconnect back
to the originating core and decrement the corresponding scratchpad entry's
pending-memory count. The entry becomes valid only after every required
response has returned. Consequently, an execution instruction can be blocked
either because a required request has not yet left the core FIFO or because
the request has left the FIFO but its response has not yet arrived. This
distinction is later used by the causal idle-time predictor.

\subsubsection{Non-Intrusive Instrumentation}

We add aggregate counters at existing state-transition points rather than
introducing new simulated pipeline stages. Per-core counters record SA-active
and SA-idle cycles, request-FIFO occupancy, and mutually exclusive execution
attribution. The interconnect records demand-active and demand-idle cycles,
and the DRAM wrapper records accepted and completed demand transactions,
channel occupancy, and transferred bytes. A memory-stage profiler accumulates
request counts and latency sums for the core queue, core-to-memory
interconnect, DRAM, and memory-to-core interconnect stages.

These counters are read-only with respect to modeled execution: they do not
change tile generation, core assignment, instruction order, operand
readiness, request routing, or completion timing. The vector-unit activity
counter is excluded from this paper because its current simulator behavior is
suspected to be incorrect. The SA, FIFO, NoC, DRAM, and request-flow counters
used in our results are retained independently of that excluded metric.
